% !TEX root = ../main.tex
\section{Introduzione all'Analisi Lineare}
L'obiettivo di questo capitolo è dimostrare la stabilità asintotica locale del velivolo VTOL nelle sue configurazioni operative critiche. La dinamica del sistema, intrinsecamente non lineare e tempo-variante, è descritta da un vettore di stato $\mathbf{x} \in \mathbb{R}^{30}$, che include posizioni, velocità, assetti angolari, ratei e gli stati integrativi dei regolatori PID \cite{simulazioneVTOL3.m}.

\section{Metodologia di Linearizzazione e Criterio di Lyapunov}
Per ogni punto di equilibrio $(\mathbf{x}_e, \mathbf{u}_e)$, il sistema è stato linearizzato tramite il calcolo della matrice Jacobiana $\mathbf{A}$:
\[ \mathbf{A} = \left. \frac{\partial \mathbf{f}(\mathbf{x}, \mathbf{u})}{\partial \mathbf{x}} \right|_{\mathbf{x}_e, \mathbf{u}_e} \]

Secondo il \textbf{Primo Metodo di Lyapunov}, la stabilità asintotica del punto di equilibrio è garantita se tutti gli autovalori $\lambda_i = \sigma_i \pm j\omega_i$ della matrice $\mathbf{A}$ presentano una parte reale $\sigma_i < 0$. Per ogni modo complesso coniugato, definiamo lo smorzamento $\zeta$ come la capacità del sistema di dissipare l'energia delle perturbazioni:
\[ \zeta = \frac{-\sigma}{\sqrt{\sigma^2 + \omega^2}} \]

\clearpage
\section{Analisi dell'Assetto di Hovering}
In fase di hovering ($V_x \approx 0$), il velivolo opera in un regime quasi-statico dove la portanza alare è nulla. Il sostentamento e il controllo del posizionamento sono affidati interamente alla spinta vettoriale dei tre rotori principali.

\subsection{Sintesi del Controllore}
Il sistema di controllo per l'hovering adotta una struttura gerarchica a doppio loop. Il loop esterno (\textit{Outer Loop}) implementa una strategia di \textbf{Sliding Mode Control (SMC)} per la regolazione della quota $z$ e delle posizioni $x, y$. L'utilizzo della funzione di commutazione $\tanh(s/\Phi)$ permette di gestire le non-linearità del modello garantendo robustezza alle perturbazioni esterne.

Il loop interno (\textit{Inner Loop}) riceve i riferimenti di assetto ($\phi_{des}, \theta_{des}$) calcolati dal controllo SMC e utilizza regolatori di tipo \textbf{PID} per stabilizzare gli angoli di Eulero. In questa fase, il momento di imbardata è controllato tramite il tilt differenziale dei rotori anteriori ($\delta_{tilt, yaw}$), mentre il rotore di coda è mantenuto all'angolo ideale $\theta_{3,ideal}$ per bilanciare la coppia residua del sistema.

\subsection{Interpretazione dei Risultati e Modi Dominanti}
L'analisi degli autovalori nel punto di equilibrio ($V_x \approx 0, z = -10$ m) conferma la stabilità asintotica locale del sistema. I risultati possono essere raggruppati secondo la separazione delle scale temporali:

\begin{itemize}
    \item \textbf{Modi ad Alta Frequenza (85.49 rad/s, $\zeta = 0.815$):} Questi autovalori sono associati alla dinamica veloce degli attuatori e ai loop interni di assetto. L'elevato smorzamento ($\zeta > 0.8$) indica che i regolatori PID sono stati correttamente tarati per prevenire risonanze strutturali e oscillazioni nei comandi di spinta.
    \item \textbf{Modi di Traiettoria (2.32 - 3.01 rad/s, $\zeta > 0.8$):} Questi modi descrivono la risposta del sistema SMC nel controllo della posizione. L'assenza di accoppiamenti aerodinamici significativi tipica del volo stazionario permette una regolazione quasi "morta" (deadbeat), dove il drone ritorna al setpoint di posizione senza sovraelongazioni apprezzabili.
    \item \textbf{Modo di Traslazione Lenta (0.44 rad/s):} Rappresenta la dinamica di deriva dovuta alla sotto-attuazione intrinseca del velivolo. Nonostante la frequenza ridotta, lo smorzamento ($\zeta = 0.578$) è sufficiente a garantire una correzione stabile della posizione orizzontale tramite il beccheggio e il rollio.
\end{itemize}

\clearpage
\section{Analisi dell'Assetto di Crociera (Cruise)}
La transizione al volo orizzontale ($V_x = 25$ m/s) sposta il sistema in un regime dominato dalle forze aerodinamiche \cite{simulazioneVTOL3.m}. Qui, la stabilità è influenzata dalla sotto-attuazione: il drone non può variare la quota senza influenzare il beccheggio e la velocità.

\subsection{Sintesi del Controllore}
Il controllo adottato è un \textit{Coordinated Bank-to-Turn}. La sfida principale è stata smorzare il modo di \textbf{Phugoid}, ovvero lo scambio energetico tra quota e velocità. Per farlo, si è agito massicciamente sui loop esterni:
\begin{itemize}
    \item Il guadagno $K_{p,x} = 12.0$ vincola rigidamente la velocità orizzontale, impedendo le oscillazioni energetiche.
    \item Il termine $K_{d,z} = 10.0$ agisce come uno smorzatore sintetico sulla velocità verticale $V_z$, "frenando" i moti oscillatori lenti.
\end{itemize}

\subsection{Discussione critica dei modi al limite}
L'analisi ha evidenziato due modi con smorzamento ridotto ($\zeta \approx 0.24$):
\begin{enumerate}
    \item \textbf{Modo Traiettoria/Dutch Roll ($2.59$ rad/s):} Rappresenta l'oscillazione accoppiata tra imbardata e rollio. Nonostante la coordinazione implementata ($K_{coord} = 0.2$), lo smorzamento rimane limitato dall'autorità dei rotori anteriori in fase di crociera.
    \item \textbf{Modo Longitudinale ($2.53$ rad/s):} È il residuo del modo di periodo corto. 
\end{enumerate}

\subsection{Vincoli di Attuazione e Trade-off dello Smorzamento}
L'analisi sperimentale condotta in regime non lineare ha evidenziato un limite fondamentale nella sintesi del guadagno derivativo $k_{d,\theta}$. Sebbene la teoria lineare suggerisca un incremento di tale parametro per migliorare lo smorzamento del modo di periodo corto, le simulazioni mostrano l'insorgere di vibrazioni ad alta frequenza sia nell'angolo di beccheggio ($\theta$) che negli angoli di tilt dei rotori anteriori.
\\Tale fenomeno è riconducibile alla natura della componente derivativa, che agisce come un amplificatore delle alte frequenze. Quando il guadagno $k_{d,\theta}$ eccede una soglia critica, il sistema di controllo tenta di rispondere a variazioni più rapide della larghezza di banda degli attuatori (fissati a circa $15$ Hz nel modello dinamico). Il risultato è un'instabilità numerica che evolve in un \textbf{ciclo limite}, dove i rotori oscillano tra i propri limiti di saturazione, degradando la stabilità globale del velivolo invece di migliorarla.
\\Lo smorzamento del modo longitudinale a valori prossimi a $\zeta = 0.25$ per l'assetto di crociera rappresenta la massima autorità di controllo esercitabile senza innescare risonanze distruttive negli attuatori.

\section{Separazione delle Scale Temporali}
L'architettura di controllo implementata si basa sul principio della separazione delle scale temporali, fondamentale per la stabilità di sistemi sotto-attuati. Tale approccio permette di decomporre il problema del controllo in due sottosistemi gerarchici:

\begin{enumerate}
    \item \textbf{Inner Loop (Dinamica Veloce):} Preposto alla stabilizzazione dei momenti e degli assetti angolari ($\phi, \theta, \psi$). I risultati della linearizzazione mostrano frequenze proprie comprese tra $60$ e $95$ rad/s, garantendo una risposta quasi istantanea alle richieste di coppia.
    \item \textbf{Outer Loop (Dinamica Lenta):} Preposto alla navigazione e al mantenimento della quota e velocità. Le frequenze dominanti in questa scala sono significativamente inferiori ($0.1 - 3$ rad/s).
\end{enumerate}

La validità di questa decomposizione è confermata dal rapporto tra le frequenze dei modi dominanti, che risulta superiore a un fattore 10 in ogni condizione di volo analizzata. Questa separazione garantisce che le correzioni di assetto avvengano a una velocità tale da poter essere considerate "stazionarie" dal punto di vista del controllo di traiettoria, minimizzando le interferenze dinamiche tra i loop e prevenendo l'insorgere di instabilità dovute all'accoppiamento dei guadagni.